
\section{Predictions and Falsification Criteria}

\subsection{Low-redshift calibration}

The Pantheon+ data release and cosmological analysis provide the supernova basis for the carried-forward tomographic calibration used here \citep{Scolnic2022,Brout2022}. The earlier scalar-topographic analysis reported
\[
A_\mu^\star=-0.0412\pm0.009
\]
in the bin
\[
0.01<z<0.03,
\]
with best-fit direction
\[
({\rm RA},{\rm Dec})
=
(211.48^\circ,-12.81^\circ).
\]
This corresponds to
\[
\frac{\delta D_L}{D_L}
=
(-0.01895\pm0.00414)\cos\theta,
\]
and at sufficiently low redshift,
\[
\frac{\delta H}{H}
=
(+0.01895\pm0.00414)\cos\theta.
\]

These numbers retain the earlier exploratory calibration; its original
likelihood is not reconstructed here. Section~\ref{sec:sn-sightlines} instead
uses the separately frozen later full-covariance amplitude
$A=-0.04321623436$ mag. The two calibrations must not be conflated, and
neither is retuned by the new interpretation test. The historical numbers
normalize only the selected physical \(n=2\) response,
\[
A_\mu
=
\frac{\mathcal K_{\rm SN}S_2}{p_2C_2},
\]
and do not independently determine \(S_2\), \(\mathcal K_{\rm SN}\), or \(C_2\). The physical distance row is now explicit; one scalar calibration leaves a state-space null direction.

\subsection{Curvature}

The pre-existing prospective curvature forecast is retained without retuning:
\[
\boxed{\Omega_k=-0.018\pm0.004}.
\]
At \(H_0=67.4\ {\rm km\,s^{-1}\,Mpc^{-1}}\),
\[
R_H
=
\frac{c}{H_0\sqrt{|\Omega_k|}}
\simeq33.15\ {\rm Gpc}.
\]
The sign prediction is \(\Omega_k<0\) in the usual FLRW curvature convention.

This value is a frozen model forecast, not a current empirical best fit. DESI DR2 BAO results are reported as well described by flat \(\Lambda\)CDM, even while combinations with other probes motivate tests of extensions such as dynamical dark energy \citep{DESIDR2ResultsII2025}. The purpose of retaining \(-0.018\pm0.004\) is therefore prospective falsification: it must be compared against independent curvature inference without being moved toward the observed result.

\subsection{Harmonic hierarchy}

The retained illustrative harmonic-filter realization predicts selective \(n=2\) softening with no analogous \(n\ge3\) enhancement. A future source-normalized \(S^3\) \(n=3\) susceptibility comparable to or larger than the physical \(n=2\) susceptibility would reject the minimal harmonic-filter branch. This statement is made at the compact-space harmonic level; comparison to a measured sky quadrupole requires an explicit observer-shell projection.

\subsection{Fixed-axis test}

A prospective independent supernova sample should first test the carried-forward axis
\[
(211.48^\circ,-12.81^\circ)
\]
rather than scanning freely over the sky. A free-axis search may be reported secondarily.

\subsection{High-n recovery}

Within the retained illustrative polynomial filter, at large \(n\),
\[
C_n\simeq\chi\ell^2n^2,
\]
so the extra scalar susceptibility falls as \(n^{-2}\). This conditional filter realization therefore predicts a strong separation between the first physical compact-space level and conventional clustering modes.

\subsection{Failure conditions}

The specified conditional realization (including its illustrative filter for the harmonic and high-$n$ tests) is disfavored if one or more of the following are established with adequate precision:
\begin{enumerate}
\item independent curvature measurements exclude the frozen negative-\(\Omega_k\) forecast;
\item the source-normalized \(S^3\) \(n=3\) response is comparable to or exceeds physical \(n=2\);
\item an independent low-redshift sample excludes the frozen dipole axis while finding a stable incompatible physical axis;
\item a coherent metric/light-cone dipole of comparable amplitude persists into a redshift regime in which the completed model predicts loss of enhanced susceptibility;
\item ordinary RSD modes show a large scale-independent modification incompatible with the high-\(n\) suppression;
\item lensing requires a metric response incompatible with the same completed scalar--metric propagator that controls matter clustering;
\item the selected physical \(n=2\) response cannot satisfy the common-domain constraint or stability requirements.
\end{enumerate}

% BEGIN R1_COSMOLOGY_CLOSURE_V15


% BEGIN PRE_PEER_REVIEW_LEGACY_STATUS_V1
\subsection{Role of carried-forward phenomenology}
The numerical low-redshift supernova calibration, fixed axis, and historical
curvature amplitude entered this project before the present geometry-first
physical reduction.  They are retained here as frozen calibration or
prospective legacy forecasts against which the new constrained construction
can be tested; they do not define the physical $n=2$ state or its observable
kernels.  In particular, the closed $S^3$ parent fixes the sign
$\Omega_k<0$ in the adopted FLRW convention, whereas the numerical forecast
$\Omega_k=-0.018\pm0.004$ is a carried-forward model forecast rather than a
new action-native determination of $|\Omega_k|$.

The same distinction applies to BAO and growth.  The current action-native
transverse distance response satisfies
$\mathcal F_{D_M/r_d}=\mathcal F_{D_L}$ conditional on
$\delta\ln r_d=0$.  The exact radial kernel
$\mathcal F_{D_H/r_d}=-\mathcal F_H$ and the survey-window/covariance
projection needed for a publishable RSD or radial-BAO comparison remain open.
The historical direct BAO--SN amplitude proportionality is therefore retained
only as provenance, not as the final prediction object.
% END PRE_PEER_REVIEW_LEGACY_STATUS_V1

\subsection{Status of the historical BAO--SN proportionality}
\label{sec:r1-bao-sn-status-v15}

The historical phenomenological forecast
$\mathcal R_{\rm BS}\sim700$--$800\,
\mathrm{km\,s^{-1}\,Mpc^{-1}}$ is retained as part of the development
record, but it is not asserted here as a closed action-native coefficient.
The modern R1 realization propagates a two-component canonical state through
observable-specific kernels, so the final falsification test is a common-state
trajectory evaluated at explicit epochs and through the appropriate standard
candle, transverse-BAO, radial-BAO, growth, or lensing response.

The low-redshift supernova calibration is kept separate from the global
high-redshift state inference. At $z_\star=0.02$ the empirical
calibration corresponds to
\[
|\Delta H_0|=1.279\pm0.279\,
\mathrm{km\,s^{-1}\,Mpc^{-1}},
\]
whereas the present high-redshift canonical state propagated to the same
low-redshift point gives a central value
\[
|\Delta H_0|_{\rm global}=0.01017\,
\mathrm{km\,s^{-1}\,Mpc^{-1}}.
\]
Consequently the minimal homogeneous/global realization is not used as a
substitute for the local calibration. Any local/topographic contribution
needed to predict that amplitude must be modeled and tested separately; no
retuning of the frozen canonical state is performed here.

The currently closed prediction structure is therefore:
\[
X_2\longrightarrow
\left\{
\mathcal F_{D_L},
\mathcal F_{D_M/r_d},
\mathcal F_{\rm growth},
\mathcal F_{\rm Weyl}
\right\},
\]
with $\mathcal F_{D_H/r_d}=-\mathcal F_H$ and survey-specific BAO
pair-window/covariance projection remaining open. These open objects are
future falsification targets rather than quantities filled by surrogate
relations.
% END R1_COSMOLOGY_CLOSURE_V15
